%============================================================
% BAB II — TINJAUAN PUSTAKA
% Struktur: A. Landasan Teori, B. Penelitian Sebelumnya.
% Bab ini sekaligus menjadi contoh pemakaian:
%   - tabel biasa (caption di atas)        : Tabel 2.1
%   - gambar  (caption di bawah)           : Gambar 2.1 & 2.2
%   - rumus matematika                     : Persamaan 2.1
%   - tabel multi-halaman "(Lanjutan)"     : Tabel 2.2
%   - tabel penelitian sebelumnya (3 artikel: 2 nasional,
%     1 internasional)                     : Tabel 2.3
%============================================================
\chapter{TINJAUAN PUSTAKA}

\section{Landasan Teori}
Menguraikan teori-teori yang menunjang penulisan/penelitian.
Referensi ataupun buku baru yang dirujuk harus dicantumkan
sebagai rujukan dan harus dicantumkan dalam daftar pustaka
\parencite{winowatan2009}. Perumusan landasan teori hendaknya
memenuhi ketentuan sebagai berikut:
\begin{enumerate}
  \item Pustaka yang digunakan terbaru, relevan, dan asli dari karya ilmiah;
  \item Apabila sumber informasi dan data yang dirujuk berasal dari buku,
        usahakan mencari terbitan edisi paling akhir;
  \item Penulisan landasan teori tidak mengandung unsur plagiat.
\end{enumerate}

Uraian teori yang memerlukan data pendukung dalam bentuk tabel
maupun gambar, format penulisan nama tabel dan gambar disajikan
seperti Tabel~\ref{tab:data-mahasiswa} dan
Gambar~\ref{fig:logo-lama}. Nomor urut tabel dan gambar
disesuaikan dengan nomor bab, diikuti dengan nomor urut dari tabel
atau gambar tersebut.

% ── Contoh tabel: caption di ATAS tabel, isi Times New Roman 10,
%    spasi 1, justify. Dipanggil di dalam paragraf (wajib). ──
\begin{tabel}{@{}c l l l@{}}
  \caption{Data mahasiswa}\label{tab:data-mahasiswa}\\
  \hline
  \rowcolor{gray!25}
  No & NIM & Nama Mahasiswa & Konsentrasi\\
  \hline
  \endfirsthead
  \TabelLanjutan{4}{Data mahasiswa}\\
  \hline
  \rowcolor{gray!25}
  No & NIM & Nama Mahasiswa & Konsentrasi\\
  \hline
  \endhead
  \hline
  \endfoot
  \endlastfoot
  1 & 11.11.0001 & Mahasiswa 1 & Konsentrasi 1\\
  2 & 11.11.0002 & Mahasiswa 2 & Konsentrasi 2\\
\end{tabel}

Terdapat perubahan bentuk logo Universitas Amikom Purwokerto,
logo sebelumnya ditunjukkan pada Gambar~\ref{fig:logo-lama},
sedangkan logo terbaru ditunjukkan pada Gambar~\ref{fig:logo-baru}.

% ── Contoh gambar: caption DI BAWAH gambar ──
\begin{figure}[h]
  \centering
  \includegraphics[width=3cm]{assets/amikom/contoh/logo_lama.png}
  \caption{Logo lama}\label{fig:logo-lama}
\end{figure}

\begin{figure}[h]
  \centering
  \includegraphics[width=4cm]{assets/amikom/contoh/logo_baru.png}
  \caption{Logo baru}\label{fig:logo-baru}
\end{figure}

Contoh penulisan rumus matematika menggunakan mode matematika
LaTeX. Tingkat akurasi suatu pengujian dapat dihitung menggunakan
persamaan~\eqref{eq:contoh-bab2}.
\begin{equation}
  Akurasi = \frac{TP + TN}{TP + TN + FP + FN} \times 100\%
  \label{eq:contoh-bab2}
\end{equation}
dengan $TP$ adalah \emph{true positive}, $TN$ adalah \emph{true
negative}, $FP$ adalah \emph{false positive}, dan $FN$ adalah
\emph{false negative}.

\section{Penelitian Sebelumnya}
Penelitian sebelumnya berisi rujukan dari jurnal atau prosiding
minimal 3 artikel yang terdiri dari 2 artikel jurnal nasional dan
1 artikel jurnal internasional \parencite{nadzif2019sistem,
haya2020rancang, anuradha2020microcontroller}. Uraikan dengan
jelas kajian pustaka yang menimbulkan gagasan dan mendasari
penelitian yang akan dilakukan, menguraikan teori, temuan, dan
bahan penelitian lain yang dijadikan landasan, serta munculkan
perbedaan Anda sebagai peneliti yang akan mengembangkan sebuah
gagasan, metode, atau hasil yang dicapai. Uraian dalam penelitian
sebelumnya diarahkan untuk menyusun kerangka atau konsep yang akan
digunakan dalam penelitian.

% ── Contoh tabel multi-halaman dengan "(Lanjutan)" ──
% Jika tabel melebihi satu halaman, longtable otomatis meneruskan
% ke halaman berikutnya dengan judul "Tabel X.Y. ... (Lanjutan)"
% dan kepala tabel yang berulang (lihat \TabelLanjutan di
% \endhead). Contoh tabel sengaja dibuat panjang agar menampilkan
% halaman lanjutan.
\begin{tabel}{@{}c l l l@{}}
  \caption{Nilai mata kuliah}\label{tab:nilai}\\
  \hline
  \rowcolor{gray!25}
  No & NIM & Nama Mahasiswa & Nilai\\
  \hline
  \endfirsthead
  \TabelLanjutan{4}{Nilai mata kuliah}\\
  \hline
  \rowcolor{gray!25}
  No & NIM & Nama Mahasiswa & Nilai\\
  \hline
  \endhead
  \hline
  \endfoot
  \endlastfoot
  1 & 11.11.0001 & Mahasiswa 1 & B\\
  2 & 11.11.0002 & Mahasiswa 2 & B\\
  3 & 11.11.0003 & Mahasiswa 3 & A\\
  4 & 11.11.0004 & Mahasiswa 4 & B\\
  5 & 11.11.0005 & Mahasiswa 5 & A\\
  6 & 11.11.0006 & Mahasiswa 6 & C\\
  7 & 11.11.0007 & Mahasiswa 7 & B\\
  8 & 11.11.0008 & Mahasiswa 8 & A\\
  9 & 11.11.0009 & Mahasiswa 9 & B\\
  10 & 11.11.0010 & Mahasiswa 10 & A\\
  11 & 11.11.0011 & Mahasiswa 11 & B\\
  12 & 11.11.0012 & Mahasiswa 12 & C\\
  13 & 11.11.0013 & Mahasiswa 13 & A\\
  14 & 11.11.0014 & Mahasiswa 14 & B\\
  15 & 11.11.0015 & Mahasiswa 15 & A\\
  16 & 11.11.0016 & Mahasiswa 16 & B\\
  17 & 11.11.0017 & Mahasiswa 17 & A\\
  18 & 11.11.0018 & Mahasiswa 18 & B\\
  19 & 11.11.0019 & Mahasiswa 19 & C\\
  20 & 11.11.0020 & Mahasiswa 20 & A\\
  21 & 11.11.0021 & Mahasiswa 21 & B\\
  22 & 11.11.0022 & Mahasiswa 22 & A\\
  23 & 11.11.0023 & Mahasiswa 23 & B\\
  24 & 11.11.0024 & Mahasiswa 24 & A\\
  25 & 11.11.0025 & Mahasiswa 25 & C\\
  26 & 11.11.0026 & Mahasiswa 26 & B\\
  27 & 11.11.0027 & Mahasiswa 27 & A\\
  28 & 11.11.0028 & Mahasiswa 28 & B\\
  29 & 11.11.0029 & Mahasiswa 29 & A\\
  30 & 11.11.0030 & Mahasiswa 30 & B\\
  31 & 11.11.0031 & Mahasiswa 31 & A\\
  32 & 11.11.0032 & Mahasiswa 32 & C\\
  33 & 11.11.0033 & Mahasiswa 33 & B\\
  34 & 11.11.0034 & Mahasiswa 34 & A\\
  35 & 11.11.0035 & Mahasiswa 35 & B\\
  36 & 11.11.0036 & Mahasiswa 36 & A\\
  37 & 11.11.0037 & Mahasiswa 37 & B\\
  38 & 11.11.0038 & Mahasiswa 38 & A\\
  39 & 11.11.0039 & Mahasiswa 39 & B\\
  40 & 11.11.0040 & Mahasiswa 40 & C\\
  41 & 11.11.0041 & Mahasiswa 41 & A\\
  42 & 11.11.0042 & Mahasiswa 42 & B\\
  43 & 11.11.0043 & Mahasiswa 43 & A\\
  44 & 11.11.0044 & Mahasiswa 44 & B\\
  45 & 11.11.0045 & Mahasiswa 45 & A\\
\end{tabel}

% ── Contoh tabel penelitian sebelumnya (minimal 3 artikel:
%    2 jurnal nasional, 1 jurnal internasional) ──
% Tabel ini dibuat LANDSCAPE (halaman miring): halaman PDF
% benar-benar dirotasi saat dilihat/dicetak, dan nomor halaman
% tetap tercetak di kanan atas dalam orientasi miring (sesuai
% contoh halaman landscape di template asli).
\begin{landskape}
\begin{tabel}{@{}c p{4.4cm} p{4.4cm} p{4.4cm} p{3.9cm} p{3.6cm}@{}}
  \caption{Penelitian sebelumnya}\label{tab:penelitian-sebelumnya}\\
  \hline
  \rowcolor{gray!25}
  No & Peneliti & Judul & Hasil Penelitian & Persamaan & Perbedaan\\
  \hline
  \endfirsthead
  \TabelLanjutan{6}{Penelitian sebelumnya}\\
  \hline
  \rowcolor{gray!25}
  No & Peneliti & Judul & Hasil Penelitian & Persamaan & Perbedaan\\
  \hline
  \endhead
  \hline
  \endfoot
  \endlastfoot
  1 & Ahmad Rudiansyah, Mardiono, dan Reza Diharja (2020) &
      Desain Alat Monitoring Kapasitas Tabung Gas LPG 3 Kilogram
      Menggunakan Load Cell Dilengkapi Dengan Deteksi Kebocoran
      Gas Berbasis Internet of Things &
      Monitoring level gas dan deteksi kebocoran gas berbasis
      Internet of Things digunakan untuk memonitoring isi gas
      LPG 3 kg dan kebocoran gas untuk memudahkan pembeli
      memantau level dan memberi peringatan ketika terjadi
      kebocoran gas. &
      Pembuatan alat monitoring level gas dan deteksi kebocoran
      gas berbasis Internet of Things &
      Menggunakan wemos D1 mini sebagai mikrokontroler;\\
  2 & Anisha Fadia Haya, Werman Kasoep, Nefy Puteri Novani (2020) &
      Rancang Bangun Smart Case Sistem Monitoring Tabung Gas
      Elpiji Berbasis Mikrokontroler &
      Mengirim data notifikasi menggunakan GSM ketika terjadi
      kebocoran gas dan memantau informasi isi dengan LCD. &
      Menggunakan load cell sensor sebagai sensor utama yang
      mendeteksi untuk monitoring level gas &
      Mikrokontroler yang digunakan adalah arduino Uno R3;\\
  3 & P. Anuradha, Rajeshwar Rao Arabelli, K. Rajkumar, dan
      J. Ravichander (2020) &
      Microcontroller Based Monitoring and Controlling of LPG
      Leaks Using Internet of Things &
      Pembuatan perangkat yang dapat mendeteksi gas menggunakan
      Arduino Uno ATmega328 dan memantau tingkat kebocoran gas
      melalui ThingSpeak. &
      Sistem untuk mendeteksi kebocoran gas menggunakan
      mikrokontroler &
      Studi ini menggunakan Arduino Uno.\\
\end{tabel}
\end{landskape}

% ── CONTOH HALAMAN LANDSCAPE (opsional) ──────────────────────
% Untuk tabel/gambar lain yang membutuhkan halaman miring,
% bungkus dengan lingkungan landskape seperti contoh di atas:
%
%   \begin{landskape}
%     \begin{tabel}{@{}c p{10cm} p{8cm}@{}}
%       \caption{Judul tabel landscape}\label{tab:landscape}\\
%       \hline
%       No & Kolom A & Kolom B\\
%       \hline
%       \endfirsthead
%       \TabelLanjutan{3}{Judul tabel landscape}\\
%       \hline
%       No & Kolom A & Kolom B\\
%       \hline
%       \endhead
%       \hline
%       \endfoot
%       \endlastfoot
%       1 & Data 1 & Data 2\\
%     \end{tabel}
%   \end{landskape}
%
% Lebar teks saat landscape ≈ 22,7 cm — sesuaikan jumlah & lebar
% kolom agar tidak melebihi batas tersebut. Nomor halaman otomatis
% tampil horizontal di kanan atas orientasi miring.
